\section{Introduction}
\label{sec:intro}

Two paradigms now dominate robot learning. \emph{Vision-language-action} (VLA)
models predict actions from frozen weights; \emph{code-as-policy} agents instead
write executable programs and improve them from experience~\citep{codeaspolicies,aspire}.
This survey is organised around that fault line.

\paragraph{Why now: the skill stack.}
The app store for robot skills has arrived; Unitree UniStore ships one-tap,
cross-model motion downloads~\citep{unistore2026}. But every skill in it is
\emph{static playback}: the least-capable point in the taxonomy. The skill stack has
three layers: \textbf{(1) Build} (how skills are created, \S\ref{sec:taxonomy}--\S\ref{sec:techniques}),
\textbf{(2) Distribute} (marketplaces and hubs), and \textbf{(3) Adapt} (on-device
self-improvement, \S\ref{sec:code}). UniStore is layer~2 over static layer-1 skills;
the promised value, ``adjust grip per fruit so nothing bruises'', lives in layer~3,
which no shipped skill does today.

\paragraph{Scope and contributions.}
This is a \emph{focused} survey rather than an exhaustive census. We examine 77 systems chosen
to populate a single analytical taxonomy, and we exclude pure perception, navigation, and
locomotion work except where it bears directly on how a manipulation skill is authored and
improved. Our contributions are: (i) a taxonomy organised by \emph{weights vs.\ skills}
(\S\ref{sec:taxonomy}); (ii) a deep-dive that arranges code-as-policy by \emph{degree of
self-improvement}, with operational definitions of the feedback, memory, and search mechanisms
and necessary-and-sufficient conditions for each rung (\S\ref{sec:code}); (iii) a map of the
``skills'' pole together with an analysis of the five senses in which ``skill'' is used
(\S\ref{sec:skill}); (iv) a summary of the characteristic strengths and limitations of each
technique family; and (v) an open-problems agenda for the emerging skill economy
(\S\ref{sec:open}). To situate these systems within the wider field,
Appendix~\ref{sec:landscape} (Table~\ref{tab:landscape}) catalogues 225 further representative
works, compared on six axes and grouped by area, bringing the survey's total coverage to more than
three hundred works.

\paragraph{Related surveys.}
Table~\ref{tab:survey_compare} contrasts the topical coverage of the closest surveys with
ours, and Figure~\ref{fig:surveys} places them on a timeline. Within \emph{ACM Computing
Surveys} the nearest are on world models, embodied intelligence, and language-model
navigation, and beyond it a cluster covers vision-language-action models and foundation models
for manipulation; but none organises code-as-policy by degree of self-improvement or names
the \S3.1.5 cell.
\input{tables/tab_survey_compare}
% Figure — main taxonomy map (§3.1–§3.6, ~47 systems), Fig2 forest style.
% Six branch colours (from styles/taxonomy_style.tex); leaves use real \citep keys.
\begin{figure*}[tbp]
\centering
\resizebox{\textwidth}{!}{%
\begin{forest}
[\textbf{Robot-learning}\\\textbf{techniques for}\\\textbf{Physical AI}, fill=tx-root, text width=11em
  [\textbf{\S3.1 Code-as-policy for robots}, fill=cat-code, for descendants={fill=leaf-code}
    [\textbf{Code-as-Policies} --- LM programs for control~\citep{codeaspolicies}]
    [\textbf{ProgPrompt} --- situated task-plan programs~\citep{progprompt}]
    [\textbf{VoxPoser} --- composable 3D value maps~\citep{voxposer}]
    [\textbf{RoboCodeX} --- multimodal behavior code~\citep{robocodex}]
    [\textbf{Instruct2Act} --- instructions to action code~\citep{instruct2act}]
    [\textbf{RoboPro} --- video-instructed policy code~\citep{robopro}]
    [\textbf{RoboCoder} --- basic skills to general tasks~\citep{robocoder}]
    [\textbf{CaP-X} --- benchmarking robot coding agents~\citep{capx}]
    [\textbf{RoboClaw} --- scalable long-horizon agent~\citep{roboclaw}]
    [\textbf{ASPIRE}~$\star$ --- agentic skill discovery + self-improvement~\citep{aspire}]
    [\textbf{ENPIRE} --- real-world policy self-improvement~\citep{enpire}]
  ]
  [\textbf{\S3.2 End-to-end VLA / generalist policies}, fill=cat-vla, for descendants={fill=leaf-vla}
    [\textbf{RT-1} --- robotics transformer at scale~\citep{rt1}]
    [\textbf{RT-2} --- web knowledge to control~\citep{rt2}]
    [\textbf{Octo} --- open generalist policy~\citep{octo}]
    [\textbf{OpenVLA} --- open-source VLA~\citep{openvla}]
    [\textbf{RoboFlamingo} --- VLM robot imitator~\citep{roboflamingo}]
    [\textbf{CogACT} --- cognition + action VLA~\citep{cogact}]
    [\textbf{SpatialVLA} --- spatial-enhanced VLA~\citep{spatialvla}]
    [$\boldsymbol{\pi_0}$ --- VLA flow model~\citep{pi0}]
    [$\boldsymbol{\pi_{0.5}}$ --- open-world generalization~\citep{pi05}]
    [\textbf{GR00T N1} --- humanoid foundation model~\citep{groot}]
    [\textbf{Gemini Robotics} --- AI in the physical world~\citep{geminirobotics}]
  ]
  [\textbf{\S3.3 LLM-authored rewards / curricula}, fill=cat-rew, for descendants={fill=leaf-rew}
    [\textbf{Eureka} --- human-level reward design~\citep{eureka}]
    [\textbf{DrEureka} --- LLM-guided sim-to-real~\citep{dreureka}]
    [\textbf{Text2Reward} --- dense reward shaping~\citep{text2reward}]
    [\textbf{Language-to-Rewards} --- reward synthesis~\citep{l2r}]
    [\textbf{Eurekaverse} --- environment curriculum gen~\citep{eurekaverse}]
    [\textbf{RoboGen} --- generative task/skill learning~\citep{robogen}]
    [\textbf{Auto MC-Reward} --- automated dense rewards~\citep{automcreward}]
  ]
  [\textbf{\S3.4 Robot skill libraries \& lifelong learning}, fill=cat-skill, for descendants={fill=leaf-skill}
    [\textbf{Voyager} --- open-ended embodied agent~\citep{voyager}]
    [\textbf{LOTUS} --- unsupervised skill discovery~\citep{lotus}]
    [\textbf{LRLL} --- lifelong robot library learning~\citep{lrll}]
    [\textbf{BOSS} --- bootstrap your own skills~\citep{boss}]
    [\textbf{SPRINT} --- scalable skill pre-training~\citep{sprint}]
    [\textbf{Uni-Skill} --- self-evolving skill repository~\citep{uniskill}]
    [\textbf{SkillFlow} --- lifelong skill discovery~\citep{skillflow}]
  ]
  [\textbf{\S3.5 Sim-to-real \& cross-embodiment transfer}, fill=cat-trans, for descendants={fill=leaf-trans}
    [\textbf{Open X-Embodiment} --- RT-X cross-embodiment~\citep{openx}]
    [\textbf{CrossFormer} --- one policy many embodiments~\citep{crossformer}]
    [\textbf{RoboCat} --- self-improving generalist agent~\citep{robocat}]
    [\textbf{Mirage} --- zero-shot cross-embodiment transfer~\citep{mirage}]
  ]
  [\textbf{\S3.6 Embodied benchmarks \& simulators}, fill=cat-bench, for descendants={fill=leaf-bench}
    [\textbf{LIBERO} --- lifelong manipulation benchmark~\citep{libero}]
    [\textbf{Robosuite} --- modular sim framework~\citep{robosuite}]
    [\textbf{BEHAVIOR-1K} --- 1000 household activities~\citep{behavior1k}]
    [\textbf{Meta-World} --- 50-task multi-task RL~\citep{metaworld}]
    [\textbf{RLBench} --- 100-task robot learning~\citep{rlbench}]
    [\textbf{ManiSkill2} --- generalizable manipulation~\citep{maniskill2}]
    [\textbf{CALVIN} --- long-horizon language tasks~\citep{calvin}]
  ]
]
\end{forest}
}
\caption{Taxonomy of robot-learning techniques for Physical AI ($\sim$47 systems).
The fault line: \S3.1 ships \emph{code} and \S3.2 ships \emph{weights}.
\textbf{ASPIRE}~$\star$ anchors the self-improving code-as-policy frontier (\S\ref{sec:code}).}
\label{fig:taxonomy_main}
\end{figure*}

% Figure -- related surveys as a compact HORIZONTAL timeline (time flows left->right).
% amber marker = published in ACM Computing Surveys; blue/teal = arXiv/MIR; red star = ours.
\begin{figure}[t]
\centering
\resizebox{\columnwidth}{!}{%
\begin{tikzpicture}[
 abv/.style={font=\scriptsize, anchor=south, align=center},
 blw/.style={font=\scriptsize, anchor=north, align=center},
 stem/.style={gray!55, line width=0.5pt},
]
 % time axis (arrow = direction of time)
 \draw[-{Stealth[length=3.2mm]}, line width=1.3pt, gray!65] (-0.4,0) -- (12.4,0);
 \node[font=\scriptsize\itshape, gray, anchor=west] at (12.1,-0.28) {time};
 % year ticks + labels (roughly proportional)
 \foreach \x in {0.3,6,10.9}{\draw[gray!35] (\x,-0.12) -- (\x,0.12);}
 \node[font=\footnotesize\bfseries,gray!80] at (1.2,-1.2) {2024};
 \node[font=\footnotesize\bfseries,gray!80] at (6.6,-1.2) {2025};
 \node[font=\footnotesize\bfseries,gray!80] at (10.1,-1.2) {2026};
 % legend
 \filldraw[fill=secD,draw=hidden-draw] (-0.35,1.42) circle (3pt);
 \node[font=\scriptsize,gray,anchor=west] at (-0.12,1.42) {= published in \emph{ACM Computing Surveys}};

 %, - events (alternating above / below), -
 \filldraw[fill=secA,draw=hidden-draw] (0.3,0) circle (3.6pt);
 \draw[stem] (0.3,0.12) -- (0.3,0.5); \node[abv] at (0.3,0.53) {\textbf{VLA survey}~\cite{survey_vla}};

 \filldraw[fill=secA,draw=hidden-draw] (2.1,0) circle (3.6pt);
 \draw[stem] (2.1,-0.12) -- (2.1,-0.5); \node[blw] at (2.1,-0.53) {\textbf{FM-Manip.}~\cite{survey_fmrobot}};

 \filldraw[fill=secD,draw=hidden-draw] (4.4,0) circle (3.6pt);
 \draw[stem] (4.4,0.12) -- (4.4,0.5); \node[abv] at (4.4,0.53) {\textbf{World Models}~\cite{survey_worldmodels}};

 \filldraw[fill=secD,draw=hidden-draw] (5.9,0) circle (3.6pt);
 \draw[stem] (5.9,-0.12) -- (5.9,-0.5); \node[blw] at (5.9,-0.53) {\textbf{Embodied Intel.}~\cite{survey_embodiedintel}};

 \filldraw[fill=secA,draw=hidden-draw] (7.4,0) circle (3.6pt);
 \draw[stem] (7.4,0.12) -- (7.4,0.5); \node[abv] at (7.4,0.53) {\textbf{Large-VLM VLA}~\cite{survey_vlmvla}};

 \filldraw[fill=secB,draw=hidden-draw] (8.9,0) circle (3.6pt);
 \draw[stem] (8.9,-0.12) -- (8.9,-0.5); \node[blw] at (8.9,-0.53) {\textbf{Embodied Learn.}~\cite{survey_embodiedlearn}};

 \filldraw[fill=secD,draw=hidden-draw] (10.1,0) circle (3.6pt);
 \draw[stem] (10.1,0.12) -- (10.1,0.5); \node[abv] at (10.1,0.53) {\textbf{Nav.-FLM}~\cite{survey_navigation}};

 % Ours -- bold red star endpoint
 \node[star,star points=5,star point ratio=2.4,fill=secE,draw=hidden-draw,minimum size=16pt,inner sep=0] at (11.7,0) {};
 \draw[secE!70!black,line width=0.6pt] (11.7,0.2) -- (11.7,0.5);
 \node[abv,text=secE!45!black] at (11.7,0.53) {\textbf{Ours}};
 \node[blw,text=secE!45!black,font=\scriptsize\itshape] at (11.7,-0.5) {+ self-improve.\\+ economy};
\end{tikzpicture}}
\caption{The closest related surveys as a timeline (2024--2026);
\textcolor{secD!55!black}{amber} markers denote surveys published in \emph{ACM Computing
Surveys} (CSUR). Prior surveys cover VLA models, foundation models, world models, and navigation,
but none organises code-as-policy by \emph{degree of self-improvement}, the weights--skills axis of
Figure~\ref{fig:corpus_collage}. \textbf{Ours}
($\star$, 2026) is organised around the self-improvement axis (\S3.1.5) and the skill economy;
per-topic coverage is in Table~\ref{tab:survey_compare}.}
\Description{A horizontal timeline of related surveys from 2024 to 2026, with amber markers for surveys
published in ACM Computing Surveys and other markers for arXiv surveys covering vision-language-action
models, foundation models, world models, and navigation; a red star at 2026 marks the present survey.}
\label{fig:surveys}
\end{figure}

\FloatBarrier
